\documentclass[12pt]{article}
\usepackage{styles/cwilliams-standard}

% Set up document metadata
\setauthors{Tzu-Chen Chiu, Arjun Sivakumar, Mara Schimmel, Christopher Williams}
\setheaderauthors{Chiu, Sivakumar, Schimmel, Williams}
\settitle{Sprint 2 Report}
\setclass{IPD}

\begin{document}

\maketitlepage

\section{Original Sprint 2 Backlog}

\begin{figure}[H]
    \centering
    \includegraphics[width=1\textwidth]{images/sprint2_backlog.png}
    \caption{Sprint 2 Backlog}
\end{figure}

\section{Sprint 2 Backlog - End of Sprint Status}

\begin{figure}[H]
    \centering
    \includegraphics[width=1\textwidth]{images/sprint2_backlog_updated.png}
    \caption{Spring 2 Backlog - Updated}
    \label{fig:sprint2updated}
\end{figure}

\section{Project Management Screenshots}

The Sprint backlogs in Google Sheets are being used to manage task.

\section{Sprint Retrospective}

\subsection*{Mara}

I initially struggled to figure out what data specifically was desired for the analog readout wrapper, but after some research I was to figure out which data 
to gather. Using pseudocode provided by our mentor, I got the measurement data, and ended up with hard and soft data for the pure detectors, pure measurements, 
and detectors gathered from the measurements. With multiple statistic data points, I believe I have accomplished my task for this sprint.

\subsection*{Chris}

This sprint I read up on and studied `Block Recurrent Transformers' (BRT) which are part of Google's Recurrent Transformer based Nueral Network. Believing that the 
I learned that the BRT is not the entirety of Google's architecture, so I did not fully complete the pipeline for this sprint and must continue work on this task.
I was able to show a working prototype of a BRT that can predict the weather temperature of an area given 48 hours of history, showcasing a proof of concept that 
we have a model that is able to learn sequence information (temperature changing over time) and that it will be useful going forward in our project.

\subsection*{Tzu-Chen}

I self-studied the structure of Transformers before I started to design the hard-input embedding. However, it is still difficult to adapt a standard Transformer, 
which is usually used for sentences or images, to the quantum-decoding use. Fortunately, there is an existing framework, PyTorch, that we can apply. It can do 
gradient descent automatically, so I don’t need to do derivatives by myself. The vectorized hard input seems to be working at this stage, so next I will study 
how to vectorize the soft input using Gaussian readout. I expect that we will need a multi-head model to handle this more complicated 2D input.

\subsection*{Arjun}

During this sprint, I completed all the tasks assigned to me for this release. The MWPM baseline error-correction work—partially built in the previous 
sprint—was finalized by cleaning and streamlining the evaluation pipeline. I also helped in developing the embedding scheme for our transformer decoder, 
collaborating closely with Tzu-Chen to resolve challenges in converting detector events and labels into vectorized, tokenized inputs suitable for model training. 
This required several coordination meetings but improved our team’s workflow and technical clarity. Additionally, I successfully connected to the ARC compute 
cluster, addressed allocation and setup hurdles, and began running model experiments on the GPUs. Overall, the sprint went smoothly with no major blockers, 
and strong teamwork helped resolve the few technical and integration issues that arose.

\section{Product Backlog (Updated)}

\begin{figure}[H]
    \centering
    \includegraphics[width=1\textwidth]{images/sprint2_updated_product_backlog.png}
    \caption{Updated Product Backlog}
\end{figure}

\begin{aside}{Picture Quality}
    We understand that this image is hard to read, but as Figure \ref{fig:sprint2updated} shows, tasks 
    3.1, 3.2, 1.7, 1.8, and 1.9 were finished while tasks 1.9 and 1.10 are still being worked on.
\end{aside}

\section{Issue Tracking}

\begin{table}[H]
\centering
\begin{tabular}{@{}llll@{}}
\toprule
\textbf{Issue} & \textbf{Description} & \textbf{Status} & \textbf{Date Completed} \\ \midrule
Analog Readout Wrapper & Configuring wrapper to output soft data & Completed & 11/21/2025 \\
AlphaQubit Structure & Learning AlpahQubit's structure & Completed & 11/17/2025 \\
Decoder & Implement MWPM into decoder code & Complete & 11/18/2025 \\
Connecting to ARC & Creating account \& connect to ARC & Completed & 11/16/2025 \\
Transformers & Learning how the heck a transformer works & In Progress & --- \\ \bottomrule
\end{tabular}
\caption{Issue Tracking Summary}
\end{table}

\section{Sprint 3 Backlog}

\begin{figure}[H]
    \centering
    \includegraphics[width=1\textwidth]{images/sprint3_backlog.png}
    \caption{Spring 3 Backlog}
\end{figure}

\section{Member Videos}

Group member videos are uploaded in the `My Media' of their Canvas accounts or linked below.

\textbf{Tzu-Chen/Jessica}: \href{https://virginiatech-my.sharepoint.com/:v:/r/personal/tzuchen_vt_edu/Documents/Attachments/TzuChen.mov?csf=1&web=1&e=YfhfV9&nav=eyJyZWZlcnJhbEluZm8iOnsicmVmZXJyYWxBcHAiOiJTdHJlYW1XZWJBcHAiLCJyZWZlcnJhbFZpZXciOiJTaGFyZURpYWxvZy1MaW5rIiwicmVmZXJyYWxBcHBQbGF0Zm9ybSI6IldlYiIsInJlZmVycmFsTW9kZSI6InZpZXcifX0%3D}{Video}



\end{document}